% -----------------------------------------------------------------------------
%  entries/qNNN.tex
%
%  One file per problem, named after its permanent ID.  The ID never changes
%  and never gets reused, even if the entry later moves to a different day.
%
%  Keys are comma-separated, so any value containing a comma must be braced:
%      topic = {Trace norm, Hilbert-Schmidt norm}
% -----------------------------------------------------------------------------
\begin{entry}{QNNN}{Short title}

\entrymeta{
  topic   = ,
  est     = $\sim$12 minutes,
  origin  = Original,          % Sourced / Extracted / Adapted / Original
  source  = Original problem,
  url     = ,                  % optional; makes "source" a live link
  prereq  = ,
  status  = drafted,           % drafted / solved / written-up / published
}

\begin{problem}
\end{problem}

\begin{solution}
\end{solution}

% \begin{remark}
% \end{remark}

\begin{takeaway}
\end{takeaway}

%  Everything below is optional.  An entry with no attempt block is complete.
%  Record only what would change what you do next time: typos and arithmetic
%  slips you caught yourself are artefacts of the medium, not facts about the
%  mathematics.  Add "public," to show this block in the public build.
\attempt{
  date       = ,               % ISO, YYYY-MM-DD -- required if the block is used
  time       = ,
  difficulty = ,               % 1-5, perceived
  howitwent  = ,
  insight    = ,
  note       = ,
}

\end{entry}
